\documentclass[../../../include/open-logic-section]{subfiles}

\begin{document}

\olfileid{sth}{choice}{banach}
\olsection{Paradoks Banach--Tarski}

Kita uga bisa nyoba mbeneraké Pilihan, kaya Boolos nyoba mbeneraké
Panggantian, kanthi panimbangan \emph{ekstrinsik} (delengen
\olref[replacement][extrinsic]{sec}). Pancen, nampa Pilihan nduwèni
akèh akibat kang dikarepaké: bisa mbandhingaké saben kardinal;
bisa ngurutaké saben himpunan kanthi becik; bisa nganggep kardinal
minangka jinis ordinal tartamtu; lan sapanunggalané.

Nanging kadhangkala ana pratelan yèn Pilihan nduwèni akibat kang
\emph{ora dikarepaké}. Umumé, pratelan iki asalé saka asil
\cite{BanachTarski1924}.

\begin{thm}[Paradoks Banach--Tarski (ing $\ZFC$)]
Sembarang bal telung dimensi bisa dipérang dadi sawatara pérangan
winates cacahé, banjur pérangan-pérangan iku bisa disusun maneh
(kanthi rotasi lan translasi) dadi rong salinan saka bal mau.
\end{thm}
\noindent
Sepisan dideleng, iki nggumunaké. Rong bal iku mesthi nduwèni volume
\emph{kaping pindho} volume bal asal. Nanging gerak kaku---rotasi
lan translasi---ora ngowahi volume. Mula katon kaya Paradoks
Banach--Tarski bisa ngetokaké matéri anyar saka ora ana.

Malah luwih nggumunaké.\footnote{Delengen
\citet[Téoréma 3.12]{Wagon2016}.} Panalaran kang padha nuduhaké yèn
bal cilik saukuran kacang polong bisa dipérang dadi sawatara pérangan
winates cacahé, banjur disusun maneh kanthi rotasi lan translasi
dadi wujud lan ukuran Big Ben.

Nanging asil-asil mau ora kabukten saka $\ZF$ waé.\footnote{Sanadyan
Banach--Tarski bisa dibuktèkaké nganggo prinsip kang luwih ringkih
tinimbang Pilihan; delengen \citet[303]{Wagon2016}.} Mula ana pilihan:
nolak Pilihan, utawa sinau nampa ``paradoks'' kasebut.

Kita nyaranaké supaya nampa ``paradoks'' kasebut. Malah, miturut kita
iki dudu paradoks kang abot. Mliginé, kita ora weruh sebabé asil iki
luwih utawa kurang paradoksal tinimbang asil-asil iki:\footnote{
\citet[276--7]{Potter2004}, \citet[16]{Weston2003}, lan
\citet[31, 308--9]{Wagon2016} ngandhakaké bab kang sarupa nganggo
conto liya.}
\begin{enumerate}
	\item Cacahé titik ing interval $(0,1)$ padha karo cacahé ing $\Real$.
	\\\emph{Bukti}: gatekna $\tan(\pi(r-\nicefrac{1}{2}))$.
	\item Cacahé titik ing garis padha karo cacahé ing persegi.
	\\Delengen \olref[his][set][pathology]{sec} lan \olref[his][set][cantorplane]{sec}.
	\item Ana kurva kang ngebaki bidang.
	\\Delengen \olref[his][set][pathology]{sec} lan \olref[his][set][hilbertcurve]{sec}.
\end{enumerate}
Telung asil iki ora mbutuhaké Pilihan. Saiki kita nganggep asil-asil
iku minangka matématika kang nggumunaké lan éndah. Mbokmenawa
Paradoks Banach--Tarski uga prayoga disawang kanthi cara kang padha.

Mesthi, ing kéné perlu ana cathetan tèknis, nanging cukup kanthi
panalaran kang ngati-ati. Gerak kaku njaga volume. Mula limang
\footnote{Kita mratelakaké paradoks iki nganggo ``sawatara pérangan
winates cacahé''. Pancen, \citet{Robinson1947} mbuktekaké yèn
pérangané bisa mung \emph{lima} (nanging ora kurang). Buktiné ana ing
\citet[kaca~66--7]{Wagon2016}.} pérangan bal mau ora bisa kabèh
\emph{bisa diukur}. Tegesé, volume ora bisa diwènèhaké marang
pérangan-pérangan iki siji-siji. Bayangna pérangan mau minangka
``sebaran tanpa wates'' saka titik-titik kang ora bisa digambar.
Mbokmenawa himpunan kang ``sumebar tanpa wates'' kaya mangkono
katon anèh. Nanging anané kaya-kaya katut saka dhawuh kang kaandhut
ing \stagesacc{}: kabèh kumpulan kang \emph{bisa} kawangun saka
himpunan kang wis ana sadurungé kudu diwujudaké.

Yèn kabèh iku durung ngyakinaké, ana panalaran ekstrinsik pungkasan
kanggo nampa Paradoks Banach--Tarski: paradoks iki langsung mènèhi
guyonan matématika kang éndah banget:
\begin{enumerate}
	\item[] \emph{Pitakonan}. Apa anagram saka ``Banach--Tarski''?
	\item[] \emph{Wangsulan}. ``Banach--Tarski Banach--Tarski''.
\end{enumerate}

\end{document}
